\documentclass[11pt,a4paper]{article}
\usepackage[T1]{fontenc}
\usepackage[utf8]{inputenc}
\usepackage{amsmath,amssymb,amsthm}
\usepackage[margin=1in]{geometry}
\usepackage[colorlinks=true,linkcolor=blue,urlcolor=blue]{hyperref}
\theoremstyle{plain}
\newtheorem{theorem}{Theorem}
\newtheorem{lemma}[theorem]{Lemma}
\newtheorem{proposition}[theorem]{Proposition}
\newtheorem{corollary}[theorem]{Corollary}
\theoremstyle{definition}
\newtheorem{remark}[theorem]{Remark}

\title{Recovering the unwritten recurrences of the table OEIS A237067}
\author{Adrian Perez Fontelles\\ \small Independent researcher}
\date{2 September 2026}
\begin{document}
\maketitle

\begin{abstract}
OEIS A237067 is a two-dimensional table $T(n,k)$ whose empirical recurrences are, for several of
its lines, given only as an ORDER --- ``k=2: [order 23]'' and the like --- with the
coefficients in a linked file rather than in the entry. Those conjectures can still be settled,
and the recurrences recovered. A line of the table is a fixed-width or fixed-height array
count, hence a walk count on finitely many vertices, hence satisfies some monic recurrence of
order at most that number; Berlekamp--Massey on twice as many exact terms returns the MINIMAL
one. Where its order equals the order the entry states --- and where the same holds after the
first few terms are dropped --- any recurrence of that order the line satisfies has a
characteristic polynomial that is a multiple of the minimal one and of equal degree, hence is
that polynomial. This note settles column 2, column 3.
\end{abstract}

\noindent\small 2020 Mathematics Subject Classification. 05A15, 11B37, 15A18.\normalsize

\section{The table and the unwritten conjectures}

OEIS A237067 is ``T(n,k)=Number of (n+1)X(k+1) 0..2 arrays with the maximum plus the upper median plus the minimum minus the lower median of every 2X2 subblock equal.''. It is read by antidiagonals and begins
\[
81, 255, 255, 826, 1101, 826, 2823, 5034,\ \dots
\]
The entry carries, among its empirical claims:
\begin{quote}\small
k=2: [order 23]\\
k=3: [order 64]
\end{quote}
As of the ``Last modified'' line on the live entry (Jul 23 2025 09:47:41, revision 6) these are still
recorded as empirical, and the recurrences themselves are not in the entry text.

\section{Why a line of the table is a walk count}

Fix the index that the line fixes. The entry defines $T(n,k)$ as a number of arrays of a shape
determined by $n$ and $k$, subject to a condition local to a bounded window of consecutive
lines. Substituting the fixed index into the entry's own wording turns the definition into an
ordinary one-parameter array count, and for such a count the admissible windows are the
vertices of a finite digraph and an array is a walk. So the line is a walk count on $S$
vertices, and by the Cayley--Hamilton theorem it satisfies a monic integer recurrence of order
at most $S$.

\section{Recovering the recurrences}

\begin{lemma}\label{lem:bm}
A sequence known to satisfy some linear recurrence of order at most $S$ is determined, with its
minimal recurrence, by its first $2S$ terms; Berlekamp--Massey applied to them returns that
minimal recurrence exactly.
\end{lemma}

\subsection*{Column $k=2$}
Substituting the fixed index into the entry's wording gives an ordinary one-parameter array
count, a walk on $S=80$ vertices. Berlekamp--Massey on $2S=160$ exact terms
returns a minimal recurrence of order $23$ --- the order the entry states --- with
integer coefficients
\begin{center}\small\begin{tabular}{cccccccc}
$c_{1}$ & $8$ & $c_{7}$ & $-59999$ & $c_{13}$ & $1512121$ & $c_{19}$ & $-395488$ \\
$c_{2}$ & $41$ & $c_{8}$ & $45046$ & $c_{14}$ & $-564301$ & $c_{20}$ & $41584$ \\
$c_{3}$ & $-446$ & $c_{9}$ & $286308$ & $c_{15}$ & $-1656226$ & $c_{21}$ & $75456$ \\
$c_{4}$ & $-123$ & $c_{10}$ & $-194469$ & $c_{16}$ & $420594$ & $c_{22}$ & $-3840$ \\
$c_{5}$ & $7323$ & $c_{11}$ & $-837761$ & $c_{17}$ & $1074464$ & $c_{23}$ & $-5760$ \\
$c_{6}$ & $-4271$ & $c_{12}$ & $441540$ & $c_{18}$ & $-181800$ &  &  \\
\end{tabular}\end{center}
so that $a(m)=\sum_{i=1}^{23}c_i\,a(m-i)$ for every $m$. Removing the first
$23+4$ terms and repeating gives order $23$ again, which is what the
identification below needs.

\subsection*{Column $k=3$}
Substituting the fixed index into the entry's wording gives an ordinary one-parameter array
count, a walk on $S=213$ vertices. Berlekamp--Massey on $2S=426$ exact terms
returns a minimal recurrence of order $64$ --- the order the entry states --- with
integer coefficients
\[
(c_1,\dots,c_{64})=(14,\ 155,\ -2926,\ -6040,\ 244521,\ -164483,\ -11263921,\ 21238672,\ 335146840,\ -838909982,\ -7072996336,\ 19721506960,\ 112388626043,\ -318347840079,\ -1397038560056,\ 3753859926320,\ 13898252323277,\ -33419242153815,\ -111947122025002,\ 228917065443925,\ 733342909262495,\ -1219480218179164,\ -3910631636706057,\ 5077215771414469,\ 16973727829173276,\ -16513175556232755,\ -59968942798729664,\ 41635759309154095,\ 172567185846981753,\ -79677765087090516,\ -404869034023921574,\ 109391191856311744,\ 775311780943181959,\ -87900710266521689,\ -1212852787383485588,\ -18159285205931616,\ 1550269674991792532,\ 186721846976722302,\ -1617927889092342392,\ -336263932844312576,\ 1375891660801136962,\ 386913316176477434,\ -949856475438318196,\ -326804386366699776,\ 529142107884198200,\ 211400822646920752,\ -235714568148881312,\ -106225582197826496,\ 82845882732048416,\ 41497774635647904,\ -22518734722941760,\ -12496183181743616,\ 4590009805465216,\ 2848783743085568,\ -666527588560896,\ -476906525168640,\ 62453272052736,\ 55743780907008,\ -2873895211008,\ -4163406077952,\ -28607201280,\ 165801000960,\ 6589513728,\ -2080899072).
\]
so that $a(m)=\sum_{i=1}^{64}c_i\,a(m-i)$ for every $m$. Removing the first
$64+4$ terms and repeating gives order $64$ again, which is what the
identification below needs.

\begin{theorem}
For each line above, the displayed recurrence holds for every index, and it is the recurrence
the entry records.
\end{theorem}

\begin{proof}
It holds by Lemma~\ref{lem:bm}. For the identification, the entry asserts that the line
satisfies SOME recurrence of the stated order, possibly only from some index on. Let $q$ be its
characteristic polynomial and $q_{\min}$ the minimal polynomial of the tail. Then
$q_{\min}\mid q$, and the second Berlekamp--Massey run --- on the line with its first few
terms removed --- shows $\deg q_{\min}$ equals the stated order, which is $\deg q$. Both are
monic, so $q=q_{\min}$.
\end{proof}

\section{Verification}

Three checks.

First, each line's model was compared against the entry's own published values for that line,
taken from the antidiagonal DATA read in either orientation and from the ``Table starts''
block, and agrees at every available term. This is the only point at which reading the entry's
English enters, and a misreading gives different counts.

Second, each recovered recurrence was evaluated directly on the model's values, with no
Berlekamp--Massey involved, and holds at every index tested.

Third, each order was computed twice by independent routes --- modulo a $61$-bit prime, where
every intermediate is a machine word, and again in exact rational arithmetic --- and once more
on the line with its first few terms dropped, which is what the identification requires. All
agree.

\begin{thebibliography}{9}
\bibitem{oeis} The OEIS Foundation, \emph{The On-Line Encyclopedia of Integer
Sequences}, \url{https://oeis.org}, 2026. Sequence A237067.
\bibitem{massey} J.~L.~Massey, Shift-register synthesis and BCH decoding, \emph{IEEE
Trans. Inform. Theory} \textbf{15} (1969), 122--127.
\bibitem{stanley} R.~P.~Stanley, \emph{Enumerative Combinatorics}, Volume 1, 2nd ed.,
Cambridge University Press, 2011. (Transfer-matrix method, Section 4.7.)
\bibitem{horn} R.~A.~Horn and C.~R.~Johnson, \emph{Matrix Analysis}, 2nd ed., Cambridge
University Press, 2013.
\end{thebibliography}

\end{document}
